\subsection{An indexed family of arithmetic theories}
\label{subsec:indexed-arithmetic-theories}

We construct one derivation object for all computably enumerable
axiom sets containing the chosen sentence axioms of PA. Put
$\mathcal B=\Tot(\Split(S_{\PA}))$, and let
$I:S_{\PA}\to\Split(S_{\PA})$ be the inclusion
$I(X)=(X,1_X)$ on objects and $I(f)=f$ for
$f:X\rightharpoonup Y$ in $S_{\PA}$.
Use the syntax realized in $\mathcal B$ by
Theorem~\ref{thm:syntactic-generated-theory}. In particular,
$S,\Fm\in\mathcal B$, and $j:S\hookrightarrow\Fm$ is the
sentence inclusion. All products, coproducts, and pullbacks in
the family construction below are taken in $\mathcal B$.
The parameter object is $E=I(\omega)=(\omega,1_\omega)$.
The interpretation $G_{\N}:\mathcal B\to\mathbf{Set}$ of
Section~\ref{sec:models} satisfies $G_{\N}E=\N$.

Choose a $\Sigma_1$ formula $\operatorname{Ax}_{\PA}(z)$
for the codes of the chosen sentence axioms of PA. Define
\[
 W(e,z)\equiv\exists y\,G(e,z,y),\qquad
 \operatorname{Ax}(e,z)\equiv
 \operatorname{Sent}(z)\land
       \bigl(\operatorname{Ax}_{\PA}(z)\lor W(e,z)\bigr).
\]
As $e$ varies over $\N$, the sets
$\{z\in\N:\N\models W(e,z)\}$ are precisely the c.e.\ subsets of
$\N$. Let $\Gamma_e$ be the set of sentences whose codes
satisfy $\operatorname{Ax}(e,z)$ in $\N$.
For each $e\in\N$, use the formula $\operatorname{Ax}(\bar e,z)$,
with $\bar e$ the numeral for $e$, to present the axiom subobject
$A_{\Gamma_e}\hookrightarrow S$. The notation
$\operatorname{Prov}_{\Gamma_e}$ below refers to this presentation.
Two enumerations with the same standard axiom codes need not define
provably equivalent predicates in $\PA$.

Let $J=J_{\PA}:\mathcal P_{\PA}\to\mathcal B$ be the
equivalence of Lemma~\ref{lem:syntactic-total-category}.
Put $\operatorname{Top}(x)\equiv(x=x)$, so that
$J(\operatorname{Top})=E$, and use the object
$S=J(\operatorname{Sent})$.
Let $\pi_0^{\mathrm{num}},\pi_1^{\mathrm{num}}:\N\to\N$
be the projections of the numerical pairing, and put
\[
 D(w)\equiv\operatorname{Sent}(\pi_1^{\mathrm{num}}w),\qquad
 A_0(w)\equiv\operatorname{Ax}
            (\pi_0^{\mathrm{num}}w,\pi_1^{\mathrm{num}}w).
\]
Let $r_0:D\to\operatorname{Top}$ and
$r_1:D\to\operatorname{Sent}$ in $\mathcal P_{\PA}$ be
represented, respectively, by
$D(w)\land e=\pi_0^{\mathrm{num}}w$ and
$D(w)\land z=\pi_1^{\mathrm{num}}w$.
The pairing identities give the product isomorphism
$v=\langle J(r_0),J(r_1)\rangle:J(D)\xrightarrow{\cong}E\times S$.
The predicate inclusion $a:A_0\hookrightarrow D$ in
$\mathcal P_{\PA}$ has graph $A_0(w)\land w'=w$.
Define the axiom subobject for the family in $\mathcal B$ by
$A=J(A_0)$ and $i=vJ(a):A\hookrightarrow E\times S$.

Let $L,R_1,R_2\in\mathcal B$ and
\[
 \lambda:L\to\Fm,\qquad
 p_n:R_n\to\Fm^n,\qquad q_n:R_n\to\Fm\quad(n=1,2)
\]
be the realized logical-axiom and rule morphisms of
Appendix~\ref{app:first-order-derivations}.
Put $F_E=E\times\Fm$ and define
\[
 R^E_0=(E\times L)+A,\qquad
 R^E_n=E\times R_n\quad(n=1,2).
\]
The nullary premise morphism is $p^E_0:R^E_0\to1$, the
unique morphism to the terminal object. Its conclusion morphism is
\[
 q^E_0=[1_E\times\lambda,(1_E\times j)i]:R^E_0\to F_E.
\]
For $n=1,2$, write
$u_n:E\times\Fm^n\to E$ and
$v_n:E\times\Fm^n\to\Fm^n$ for the product projections,
and $\pi_{n,k}:\Fm^n\to\Fm$ for the $k$th projection,
where $1\le k\le n$.
Define $d_n:E\times\Fm^n\to F_E^n$ by its components
$\langle u_n,\pi_{n,k}v_n\rangle:E\times\Fm^n\to F_E$,
and put
\[
 p^E_n=d_n(1_E\times p_n):R_n^E\to F_E^n,\qquad
 q^E_n=1_E\times q_n:R_n^E\to F_E.
\]
Every premise and the conclusion of a rule have the same
$E$-coordinate. These morphisms define the inference endofunctor
$H_E:\mathcal B/F_E\to\mathcal B/F_E$
of Appendix~\ref{app:closed-subobjects}.
Lemma~\ref{lem:syntactic-derivation-initiality} gives its
initial algebra
\[
 \mathfrak P=((P^{\mathrm{all}},c^{\mathrm{all}}),\alpha),
 \qquad c^{\mathrm{all}}:P^{\mathrm{all}}\to F_E,
\]
where
$\alpha:H_E(P^{\mathrm{all}},c^{\mathrm{all}})
             \to(P^{\mathrm{all}},c^{\mathrm{all}})$
is a morphism in $\mathcal B/F_E$.
Let $c:P\to E\times S$ be the pullback of $c^{\mathrm{all}}$
along $1_E\times j:E\times S\to F_E$, and let
$m:\mathcal T\hookrightarrow E\times S$ be the regular image
of $c$. Theorem~\ref{thm:least-theorem-subobject} identifies
the image of $c^{\mathrm{all}}$ as the least subobject closed
under these rules. Stability of regular images identifies
$m$ with the pullback of that image along $1_E\times j$.
By Lemma~\ref{lem:syntactic-total-category}, $m$ has a partial
inverse $r:E\times S\rightharpoonup\mathcal T$ in
$\Split(S_{\PA})$. Thus $rm=1_{\mathcal T}$ and $mr=\bar r$;
Lemma~\ref{lem:range-regular-image} gives
\[
 t=mr=\ran(c):E\times S\rightharpoonup E\times S
 \quad\text{in }\Split(S_{\PA}).
\]

For $e\in\N$, let $\epsilon_e:1\to E$ be the numeral
morphism. For $X\in\mathcal B$, let $!_X:X\to1$ be the
terminal morphism, and define
$k_e^X=\langle\epsilon_e!_X,1_X\rangle:X\to E\times X$.

\begin{proposition}[Fibers of the derivation family]
\label{prop:arithmetic-theory-fibers}
For every $e\in\N$, the pullback of $c^{\mathrm{all}}$ along
$k_e^{\Fm}$ is the underlying conclusion morphism of an initial
$H_{A_{\Gamma_e}}$-algebra in $\mathcal B/\Fm$.
There is an induced morphism $\iota_e:P_{\Gamma_e}\to P$
in $\mathcal B$, and the square
\[
\begin{tikzcd}
 P_{\Gamma_e}\arrow[r,"\iota_e"]
       \arrow[d,"c_{\Gamma_e}"'] & P\arrow[d,"c"]\\
 S\arrow[r,"k_e^S"'] & E\times S
\end{tikzcd}
\]
is a pullback in $\mathcal B$. Consequently
$(k_e^S)^*m$ is the sentence theory generated by $A_{\Gamma_e}$.
\end{proposition}
\begin{proof}
Put $E_{\ne e}=(\omega,[x\ne\bar e]_{\PA})\in\mathcal B$,
and let $u_e:E_{\ne e}\hookrightarrow E$ be its inclusion.
The predicates $x=\bar e$ and $x\ne\bar e$ give the
isomorphism $[\epsilon_e,u_e]:1+E_{\ne e}\to E$.
For $X\in\mathcal B$, put
$k_{\ne e}^X=u_e\times1_X:E_{\ne e}\times X\to E\times X$.
Finite extensivity gives an equivalence
\[
 \Xi_e:\mathcal B/F_E\longrightarrow
       \mathcal B/\Fm\times\mathcal B/(E_{\ne e}\times\Fm).
\]
For an object $(Q,q)$ of $\mathcal B/F_E$, define
\[
 \Xi_e(Q,q)=((Q_e,q_e),(Q_{\ne e},q_{\ne e}))
   =\bigl((k_e^{\Fm})^*(Q,q),(k_{\ne e}^{\Fm})^*(Q,q)\bigr).
\]
Thus $q_e:Q_e\to\Fm$ and
$q_{\ne e}:Q_{\ne e}\to E_{\ne e}\times\Fm$ are pullbacks
of $q$. For a slice morphism $f:(Q,q)\to(Q',q')$, represented
by $f:Q\to Q'$ with $q'f=q$, put
\[
 \Xi_e(f)=(f_e,f_{\ne e}),\qquad
 f_e:(Q_e,q_e)\to(Q'_e,q'_e),\quad
 f_{\ne e}:(Q_{\ne e},q_{\ne e})\to(Q'_{\ne e},q'_{\ne e}),
\]
where each component is the induced pullback morphism in its
respective slice category. Pullback uniqueness proves the
identity and composition laws.
A quasi-inverse
\[
 \Psi_e:\mathcal B/\Fm\times\mathcal B/(E_{\ne e}\times\Fm)
                 \longrightarrow\mathcal B/F_E
\]
sends $(x:X\to\Fm,y:Y\to E_{\ne e}\times\Fm)$ to
$[k_e^{\Fm}x,k_{\ne e}^{\Fm}y]:X+Y\to F_E$.
For $f:(X,x)\to(X',x')$ in $\mathcal B/\Fm$ and
$g:(Y,y)\to(Y',y')$ in $\mathcal B/(E_{\ne e}\times\Fm)$,
its action on morphisms is represented by
$f+g:X+Y\to X'+Y'$ in $\mathcal B$.
The decomposition $Q\cong Q_e+Q_{\ne e}$ supplies the
equivalence isomorphisms.

Write $F_{\ne e}=E_{\ne e}\times\Fm$, and let
$\pi_E:F_E\to E$ be the product projection. For $n=0,1,2$,
put $\rho_n=\pi_Eq_n^E:R_n^E\to E$, and form the pullback
\[
 R_n^{\ne e}=E_{\ne e}\times_E R_n^E,
 \qquad s_n:R_n^{\ne e}\to R_n^E,
 \qquad t_n:R_n^{\ne e}\to E_{\ne e},
 \qquad \rho_ns_n=u_et_n.
\]
Since every premise has the parameter of the conclusion, there
are unique morphisms
\[
 p_n^{\ne e}:R_n^{\ne e}\to F_{\ne e}^n,
 \qquad q_n^{\ne e}:R_n^{\ne e}\to F_{\ne e}
\]
satisfying
\[
 (k_{\ne e}^{\Fm})^n p_n^{\ne e}=p_n^Es_n,
 \qquad k_{\ne e}^{\Fm}q_n^{\ne e}=q_n^Es_n.
\]
These are the premise and conclusion morphisms for an inference
endofunctor
$H_{\ne e}:\mathcal B/F_{\ne e}\to\mathcal B/F_{\ne e}$.
Write $H_e=H_{A_{\Gamma_e}}:\mathcal B/\Fm\to\mathcal B/\Fm$
for the inference endofunctor of $A_{\Gamma_e}$.
There is a natural isomorphism
\[
 \zeta:\Xi_eH_E\xRightarrow{\cong}
                 (H_e\times H_{\ne e})\Xi_e.
\]
Its component at $(Q,q)$ consists of two slice morphisms
\[
 \begin{aligned}
 \zeta^0_{(Q,q)}&:(k_e^{\Fm})^*H_E(Q,q)\xrightarrow{\cong}
                                      H_e(Q_e,q_e),\\
 \zeta^1_{(Q,q)}&:(k_{\ne e}^{\Fm})^*H_E(Q,q)\xrightarrow{\cong}
                                      H_{\ne e}(Q_{\ne e},q_{\ne e}).
 \end{aligned}
\]
The first belongs to $\mathcal B/\Fm$ and the second to
$\mathcal B/F_{\ne e}$.
For the nullary summand of $H_E(Q,q)$, the two pullbacks of
$R_0^E$ are the corresponding axiom objects. In the first,
the nonlogical axiom predicate is $\operatorname{Ax}(\bar e,z)$.
For $n=1,2$, the equation $p_n^E=d_n(1_E\times p_n)$ makes
the parameter of every premise equal to that of the rule.
Pullback pasting therefore gives the displayed isomorphisms on
each summand of arity $n$. Taking their coproducts defines
$\zeta^0_{(Q,q)}$ and $\zeta^1_{(Q,q)}$.
For $f:(Q,q)\to(Q',q')$ in $\mathcal B/F_E$, uniqueness
of the pullback morphisms gives the naturality equation
\[
 (H_e\times H_{\ne e})(\Xi_e f)\,\zeta_{(Q,q)}
   =\zeta_{(Q',q')}\,\Xi_e(H_Ef).
\]
Both sides are morphisms from $\Xi_eH_E(Q,q)$ to
$(H_e\times H_{\ne e})\Xi_e(Q',q')$ in the product of
the two slice categories.

The equivalence $\Xi_e$ and the natural isomorphism $\zeta$
induce an equivalence
\[
 \Lambda_e:\operatorname{Alg}(H_E)\longrightarrow
             \operatorname{Alg}(H_e)\times
             \operatorname{Alg}(H_{\ne e}).
\]
An algebra $((Q,q),\delta)$, where
$\delta:H_E(Q,q)\to(Q,q)$ in $\mathcal B/F_E$, is sent
to the pair of algebras whose structure morphisms are the
components of
\[
 \Xi_e(\delta)\,\zeta_{(Q,q)}^{-1}:
 (H_e\times H_{\ne e})\Xi_e(Q,q)\longrightarrow\Xi_e(Q,q).
\]
For an algebra morphism
$a:((Q,q),\delta)\to((Q',q'),\delta')$ in
$\operatorname{Alg}(H_E)$, let
$f:(Q,q)\to(Q',q')$ be its underlying slice morphism.
Then $\Lambda_e(a)$ has underlying slice morphisms
$(f_e,f_{\ne e})$. Naturality of $\zeta$ identifies the
algebra-morphism equations in the two categories. Full faithfulness
of $\Xi_e$ therefore gives full faithfulness of $\Lambda_e$.
Transporting a pair of structure morphisms through $\Psi_e$
gives essential surjectivity.

Let $(\mathfrak P_e,\mathfrak P_{\ne e})=\Lambda_e(\mathfrak P)$.
This pair is initial in
$\operatorname{Alg}(H_e)\times\operatorname{Alg}(H_{\ne e})$.
For any $\mathfrak Q\in\operatorname{Alg}(H_e)$, there is
therefore a unique morphism
\[
 (\mathfrak P_e,\mathfrak P_{\ne e})\longrightarrow
 (\mathfrak Q,\mathfrak P_{\ne e})
\]
in that product category. Its first component is an
$H_e$-algebra morphism $\mathfrak P_e\to\mathfrak Q$.
Pairing any two such algebra morphisms with
$1_{\mathfrak P_{\ne e}}$ proves their equality. Hence
$\mathfrak P_e$ is initial and is therefore isomorphic in
$\operatorname{Alg}(H_e)$ to the initial derivation algebra for
$A_{\Gamma_e}$ from Theorem~\ref{thm:syntactic-generated-theory}.
The underlying slice object of $\mathfrak P_e$ is
$(k_e^{\Fm})^*(P^{\mathrm{all}},c^{\mathrm{all}})$.

Pull the underlying slice isomorphism of the initial algebras back along
$j:S\hookrightarrow\Fm$. Since $(1_E\times j)k_e^S=k_e^{\Fm}j$,
pullback pasting gives an isomorphism
$P_{\Gamma_e}\cong S\times_{E\times S}P$ over $S$.
Composing it with the projection to $P$ gives $\iota_e$ in the
displayed square. Regular-image stability shows that
$(k_e^S)^*m$ is the image of $c_{\Gamma_e}$.
Appendix~\ref{app:first-order-derivations} identifies that image
with the generated sentence theory.
\end{proof}

Choose a $\Sigma_1$ predicate $P_0(d)$ and an
isomorphism $u:J(P_0)\xrightarrow{\cong}P$ in $\mathcal B$.
By fullness of $J$, there is a morphism
$\gamma_c:P_0\to D$ in $\mathcal P_{\PA}$ satisfying
\[
 J(\gamma_c)=v^{-1}cu:J(P_0)\longrightarrow J(D).
\]
Choose a standard program index $a_c\in\N$ whose graph
$G_{a_c}(d,w)$ represents $\gamma_c$, as in
Lemma~\ref{lem:syntactic-functional-search}. The defining conditions
\eqref{eq:syntactic-support} and \eqref{eq:syntactic-totality}
for $\gamma_c:P_0\to D$ give
\[
 \PA\vdash\forall d\,
     \bigl(\exists w\,G_{a_c}(d,w)\leftrightarrow P_0(d)\bigr).
\]
Define $\operatorname{Prov}(e,z)
\equiv\exists d\,G_{a_c}(d,\langle e,z\rangle)$.
The product isomorphism $v$ makes this a $\Sigma_1$
predicate for $m:\mathcal T\hookrightarrow E\times S$.
For each $e\in\N$,
Proposition~\ref{prop:arithmetic-theory-fibers} identifies the
represented subobjects of $S$, giving
\[
 \PA\vdash\forall z\,
   \bigl(\operatorname{Prov}(\bar e,z)
       \leftrightarrow\operatorname{Prov}_{\Gamma_e}(z)\bigr).
\]

